% Chapter 2: Problem definition and scope

\chapter{Problem Definition and Scope}
\label{ch:problem-statement}

\section{Faithfulness as a behavioural question}
\label{sec:problem-definition}

An explanation is faithful when it accurately represents the reasoning process
that produced a model prediction \citep{DBLP:conf/acl/JacoviG20}. This
definition concerns a causal relation inside the model. Generated text alone
does not reveal that relation. A rationale can be fluent, relevant, and
consistent with an answer while omitting the evidence that changed the answer.

We therefore ask a narrower, observable question. If an answer or rationale
depends on visual evidence, controlled changes to that evidence should sometimes
change the corresponding output. Stage~1 measures answer
dependence. Stage~2 measures how the rationale text responds while the supplied
answer is held fixed. These results describe behavioural dependence on the
input. They do not reveal the hidden computation that produced an output.

This distinction also separates two types of stability. A stable output under
horizontal mirror can be compatible with a transformation that preserves scene
content. The same stability under grey replacement or an unrelated image has a
different interpretation because content relevant to the answer has been
removed or replaced. The six image conditions make these comparisons explicit.

\section{Answers and rationales are separate outcomes}
\label{sec:problem-two-stage}

The two-stage design prevents answer accuracy from standing in for rationale
quality. Stage~1 maps the image, question, choices, and optional description to
an answer. Stage~2 maps the same record and a supplied answer to a rationale.
The two stages use separate adaptations of the same base model for each arm.

The two stages are related but distinct. Stage~2 begins from the arm's trained
Stage~1 weights and is then trained further to write a rationale for an answer
it is given. It does not receive Stage~1's hidden states, its four class scores,
or any trace of how that answer was reached. The classifier head belongs to
Stage~1 alone. A Stage~2 rationale therefore describes the answer it was handed,
and the design tests whether each output depends on the visual evidence when the
same intervention is applied to both.

Gold answer conditioning is central to this separation: it keeps the answer
that Stage~2 must support constant while the image changes
(Section~\ref{sec:setup-regimes}). The primary RQ2 analysis holds the answer
fixed in this way, while a separate source-only regime examines the recorded
path in which Stage~2 receives a Stage~1 prediction instead
(Section~\ref{sec:setup-pred-a}).

Accuracy and Drift use different scales and denominators. Their relationship is
therefore assessed through their patterns across image conditions rather than by
comparing their numerical magnitudes.

\section{Input arms and image conditions}
\label{sec:problem-design}

The four input arms vary two inputs. The \texttt{point} arms add polygon
overlays to the image. The \texttt{desc} arms add a fixed textual description.
Each arm has its own Stage~1 and Stage~2 checkpoint, so a comparison between
arms describes what these four adapted systems do, rather than the effect of
one input component on its own.

The six image conditions vary the evidence available during evaluation. Source
provides the image appropriate to the arm. Grey removes scene content. Mismatch
supplies an unrelated image from a different source frame. Mask and noise
corrupt the image while retaining part of its structure. Mirror preserves the
scene but changes its spatial orientation. Chapter~\ref{ch:setup} defines each
construction.

\section{Scope and interpretation}
\label{sec:problem-scope}

The evaluation covers one Qwen2.5-VL model size, VCR, one training seed per
arm, one prompt family, and deterministic generation. Within each stage and
arm, all image conditions use the same checkpoint. Across arms, the checkpoints
were trained separately, so the observed differences may reflect both the input
configuration and variation introduced during training. We therefore interpret
comparisons between arms descriptively, as properties of the four trained
systems.

The description arms retain text derived from the image after the pixels are
changed, so smaller accuracy losses or lower Drift in those arms can reflect
textual substitution. A polygon overlay changes the pixels and a description
adds text, so the analysis reports their patterns separately.

Every estimate is also conditional on the eight checkpoints the study used.
Each arm was adapted once. Stage~1 selected its checkpoint by validation
accuracy on the same records that the results report, and Stage~2 retained the
final checkpoint of a fixed epoch budget scored on those records
(Section~\ref{sec:setup-model}). The intervals in Chapter~\ref{ch:results}
describe variation across records and source frames for these fixed
checkpoints.
